\documentclass[conference]{IEEEtran}
\usepackage[utf8]{inputenc}
\usepackage[T1]{fontenc}
\usepackage{cite}
\usepackage{amsmath,amssymb}
\usepackage{graphicx}
\usepackage{booktabs}
\usepackage{url}

\title{Does Deterministic Generate--Validate--Repair Reduce Unsafe LLM‑Generated Network Changes? A Confirmatory Paired Study}

\author{
\IEEEauthorblockN{Autonomous AI Researcher}
\IEEEauthorblockA{CNSM 2026 Agentic AI Researcher Track}
}

\begin{document}

\maketitle

\begin{abstract}
We preregistered and executed a sealed paired confirmatory study (n = 200 intent\ensuremath{\rightarrow}configuration tasks) to test whether a deterministic generate--validate--repair (GVR) pipeline that consults a deterministic NetOps validator reduces validator‑detected unsafe or semantically incorrect network changes relative to an LLM‑only generation pipeline. The run used a single hosted model (gpt‑5‑mini) with adapter‑enforced shared‑initial generation semantics and a primary deterministic validator (deterministic\_netops\_validator\_v5). Execution completed all planned pairs (completed\_pairs = 200) consuming 200 model calls (one shared initial generation per task); no repair calls were issued because every initial candidate passed the validator. Paired contingency: n11 = 200, n10 = 0, n01 = 0, n00 = 0; estimate = 0.0; exact McNemar two‑sided p = 1.0; 95\% bootstrap percentile CI = [0.0, 0.0] (10,000 resamples, seed = 1729). Because the baseline validator‑detected unsafe incidence was 0\% on this task bank, the preregistered >=50\% relative‑reduction hypothesis could not be meaningfully evaluated. We report the sealed design, execution accounting (start: 2026‑08‑30T11:50:34.530059Z; end: 2026‑08‑30T12:12:07.541650Z), the Mininet 10\% holdout (20 tasks) that produced zero validator--emulator disagreements, and an artifact‑grounded discussion of operational implications and threats to validity (preregistration id: cnsms2026‑gvr‑effectiveness‑02).
\end{abstract}

\section{Introduction}
Automating intent\ensuremath{\rightarrow}configuration translation with generative language models can increase NetOps productivity but risks producing incorrect or unsafe changes that cause outages or policy violations. Prior work therefore proposes grounding, deterministic validation, and repair loops (generate--validate--repair, GVR) to detect and correct unsafe candidates before application. The operational benefit of automated repair is conditional on three factors: (i) baseline prevalence of validator‑detectable failures from the chosen generator/prompting policy; (ii) validator coverage and fidelity; and (iii) the available repair budget and repair‑prompt effectiveness. We preregistered a narrowly scoped confirmatory question: Does an automated, deterministic GVR pipeline that consults a deterministic NetOps validator materially reduce validator‑detected unsafe or semantically incorrect network changes relative to an LLM‑only generation pipeline on a curated 200‑task intent\ensuremath{\rightarrow}configuration bank? The study was designed as a paired experiment that fixes the generator output per task and isolates the marginal effect of invoking a validator+repair on the same candidate.

\section{Related Work}
Surveys and empirical studies across LLM applications document factual errors, hallucinations, prompt sensitivity, and non‑determinism; mitigation approaches include retrieval grounding, verifier‑assisted repair, and constrained agent architectures. NetOps prototypes show feasibility of intent\ensuremath{\rightarrow}configuration generation but typically lack instrumented, production‑like failure‑rate evaluations. Architectural proposals recommend combining LLM generation with deterministic validators (e.g., Batfish‑style checkers) and with guarded agent patterns; red‑team audits demonstrate configuration‑level brittleness and tool‑description injection vectors that can alter safety outcomes. Our study positions itself as an externally auditable, preregistered measurement that isolates the marginal effect of deterministic repair for a fixed model and sampling policy, addressing gaps about when GVR actually reduces validator‑detected failures in practice.

\section{Methodology}
Preregistration and inferential plan. We froze the experimental plan under preregistration id cnsms2026‑gvr‑effectiveness‑02 before any model calls. The confirmatory design was paired and deterministic: n = 200 tasks (four topology clusters \ensuremath{\times} 50 tasks), one shared deterministic initial LLM generation per task, and a repair budget of at most one automated repair call per task. The primary estimand was the paired\_success\_rate\_difference\_guarded\_minus\_baseline and the prespecified test was an exact two‑sided McNemar test (\ensuremath{\alpha} = 0.05). Confidence intervals for the paired difference were computed by the bootstrap percentile method with 10,000 resamples and bootstrap\_seed = 1729. The analysis executor was paired\_binary\_analysis\_v1 as recorded in the analysis artifacts.

Adapter, model, and sampling semantics. The run used a single hosted model declared as gpt‑5‑mini (provider label: openai\_responses). Model configuration metadata (execution/model\_configuration.json) records requested\_model = "gpt-5-mini", max\_output\_tokens = 2000, maximum\_attempts\_per\_call = 1, and temperature = null (deterministic sampling enforced by the adapter). Crucially, the hosted\_netops\_gvr\_v1 adapter enforced shared\_initial generation semantics: for each task the system generated exactly one initial candidate which was deterministically reused for both the baseline (LLM‑only) and guarded (GVR) conditions. This shared‑initial design eliminates per‑condition sampling variance and isolates the incremental effect of invoking validation and an allowed repair step on the identical generator output. Adapter accounting in deterministic\_reconciliation.json documents hosted\_model\_call\_event\_count = 200, stage\_counts = \{shared\_initial\_generation: 200\}, and model\_calls\_used = 200, confirming the enforced single shared invocation per task and that no separate condition‑specific sampling occurred.

Task bank, validators, and holdout emulator. The task bank is a curated, machine‑checkable intent\ensuremath{\rightarrow}configuration set (200 tasks) in which each task encodes initial and expected final states plus workflow policies (e.g., make\_before\_break\_uplink\_switchover, must\_be\_down\_for\_fields). The primary deterministic checker (deterministic\_netops\_validator\_v5, validator\_version = "5.0") runs syntactic parsing, intent/constraint enforcement, and structured diagnostics (fields available in validator traces include valid flag, violation\_count, violation\_codes, parsed\_command\_count, required\_command\_count, and difficulty annotations). A Mininet‑based data‑plane emulator was preregistered as a secondary disagreement detector applied to a randomized 10\% holdout (20 tasks) to detect validator--emulator mismatches; the executed contamination\_summary artifact reports zero disagreement flags for that holdout. The preregistration specified deterministic treatment rules: pairs flagged for validator--emulator disagreement would be excluded from the primary analysis per the contamination\_plan; sensitivity analyses for alternative treatments were preregistered but not invoked because no flags occurred.

Repair policy, budgeting, and missingness handling. Per the preregistration, when the primary validator reported violations the pipeline was permitted to issue up to one automated repair prompt to the same LLM for that task (maximum repair calls across the sample = 200). Repair prompts, if used, would be captured in repair\_provider\_trace fields in the episode artifacts; none appear in the archived episodes because no repairs were applied. Missingness and exclusion rules were deterministic: pairs where neither condition completed because of infrastructure failures would be excluded and logged; pairs with contamination flags (validator--emulator disagreement) would be excluded from the primary confirmatory test. The execution manifest and missingness\_summary show complete\_pair\_count = 200 and zero failed episodes, so the planned missingness handling was not invoked in practice.

Execution accounting and artifact capture. The autonomous run executed between 2026‑08‑30T11:50:34.530059Z and 2026‑08‑30T12:12:07.541650Z. The adapter enforced the preregistered maximum model calls (maximum\_model\_calls = 400) and per‑task caps (initial generation = 1 call; repair calls <= 1). The run consumed 200 model calls (one shared initial generation per task), with cache\_miss\_count = 200 and cache\_hit\_count = 0 as recorded in the deterministic reconciliation artifacts. All model calls, prompts, raw responses, validator traces, provider traces, and analysis artifacts (execution/raw\_results.jsonl, analysis/paired\_contingency\_table.csv, analysis/condition\_summary.csv, analysis/contamination\_summary.csv, analysis/analysis\_log.jsonl, execution/model\_configuration.json) were archived; representative per‑task artifacts (e.g., execution/scoring/task-000001-baseline.json) show validator\_trace fields and scoring outcomes. The analysis pipeline executed the exact McNemar test and bootstrap CI generation exactly as preregistered; analysis/analysis\_log.jsonl records bootstrap\_resamples = 10000 and bootstrap\_seed = 1729.

Design rationale and implications for internal validity. The shared\_initial paired design was chosen to isolate the pure marginal effect of validation+repair on a fixed generator output; this reduces variance attributable to stochastic sampling and makes the confirmatory test focused on whether repair can convert a validator‑failing candidate into a validator‑passing one for the same initial output. The trade‑off is explicit: the design sacrifices generality across multiple independent generator draws for stronger internal control. Preregistration, deterministic adapter semantics, and frozen stopping rules were used to avoid post‑hoc analytic choices. Deterministic reconciliation artifacts confirm pair\_accounting\_consistent = true and all\_deterministic\_consistency\_checks\_passed = true.

Limitations baked into the methodology. The methodology intentionally restricts scope: a single model, a deterministic prompt/sampling policy, a one‑repair budget, and a curated task bank. These constraints reflect an explicit confirmatory question about marginal GVR effect under a tightly controlled generator policy rather than an evaluation of alternative sampling, multi‑step repair strategies, or cross‑model generality. When interpreting outcomes, readers should account for these preregistered methodological choices documented in the archived plan.

\section{Results}
Execution completeness and basic counts. The experiment executed to completion under the preregistered stopping rule. Accounting artifacts report completed\_episode\_count = 400 (shared initial generation reused across conditions), complete\_pair\_count = 200, failed\_episode\_count = 0, and model\_calls\_used = 200 (one shared initial generation per task). The adapter/provider audit recorded hosted\_model\_call\_event\_count = 200, cache\_miss\_count = 200, stage\_counts = \{shared\_initial\_generation: 200\}, and condition\_counts = \{shared: 200\}, confirming that no separate per‑condition re‑sampling occurred.

Validator outcomes and paired contingency. The primary deterministic validator (deterministic\_netops\_validator\_v5) marked every shared initial candidate as valid. The condition summary artifact (analysis/condition\_summary.csv) contains: "baseline,200,0,200,1.0" and "guarded,200,0,200,1.0" (success\_rate = 1.0 for both conditions). The paired contingency artifact (analysis/paired\_contingency\_table.csv) records the contingency table used in the exact McNemar test as:
- n11 (pass/pass) = 200
- n10 (guarded pass, baseline fail) = 0
- n01 (guarded fail, baseline pass) = 0
- n00 (fail/fail) = 0

Inferential result and bootstrap CI. With zero discordant pairs the observed paired difference estimate = 0.0. The pre‑registered exact McNemar two‑sided test returns p = 1.0. The bootstrap percentile 95\% confidence interval computed with 10,000 resamples and bootstrap\_seed = 1729 is [0.0, 0.0]. These numerical outputs are recorded in analysis/results.json and the analysis/analysis\_log.jsonl artifact (bootstrap\_resamples = 10000, bootstrap\_seed = 1729).

Repair activity and secondary outcomes. The preregistered repair budget allowed up to 200 repair calls (one per task) but no repair calls were issued because no initial candidate failed the validator. The secondary\_results array is therefore empty and the intended secondary estimands (average repair iterations, GVR‑introduced failure rate, model‑call cost per successful repair) are unpopulated in the analysis bundle. The analysis logs explicitly note repair\_calls\_used = 0.

Representative per‑task diagnostics. Representative scoring artifacts for sample pairs illustrate validator traces and complexity indicators while still producing validator‑passing configurations. Examples taken from archived representative tasks:
- pair‑000001: validator\_trace.validation\_after.valid = true; violation\_count = 0; parsed\_command\_count = 4; required\_command\_count = 4; difficulty.level = "medium"; pattern = "shutdown\_configure\_restore"; required\_sequence length = 4.
- pair‑000002: validator\_trace.validation\_after.valid = true; violation\_count = 0; parsed\_command\_count = 2; required\_command\_count = 2; difficulty.level = "hard"; pattern = "make\_before\_break\_uplink\_switchover".
- pair‑000003: validator\_trace.validation\_after.valid = true; violation\_count = 0; parsed\_command\_count = 6; required\_command\_count = 6; difficulty.level = "hard"; pattern = "paired\_link\_atomic\_mtu\_change".
These artifacts demonstrate that tasks exercised medium‑to‑very\_hard workflow constraints (examples in the archive include labels "medium", "hard", and "very\_hard") and that the model outputs conformed to the task‑encoded required\_sequences and validator checks despite nontrivial required\_command\_counts.

Emulator (contamination) check. Per the preregistered contamination\_plan, a Mininet‑based data‑plane emulator was run on a randomized 10\% holdout (20 tasks). The analysis/contamination\_summary.csv artifact records zero disagreement flags between deterministic\_netops\_validator\_v5 and the Mininet emulator on that holdout (flag counts empty). Deterministic reconciliation artifacts further report pair\_accounting\_consistent = true and all deterministic consistency checks passed. We note, however, that the explicit labeled seed for the randomized holdout selection is not present among the labeled artifacts; the contamination\_summary artifact nevertheless documents that the secondary emulator check executed and found no disagreements for the recorded holdout.

Unexecuted sensitivity branches and their artifact evidence. Several preregistered sensitivity analyses and contingency branches (e.g., treating flagged tasks as failures or successes, alternative missingness treatments, and inclusion of flagged tasks in the primary analysis) were not invoked because no tasks were flagged and no episodes failed. The execution manifest and analysis artifacts explicitly record these branches as unused; the analysis/contamination\_summary.csv and analysis/missingness\_summary.csv reflect zero flags and zero missing pairs.

Operationally relevant diagnostics. The run produces three operationally meaningful, artifact‑grounded diagnostics: (1) baseline validator failure prevalence in this curated workload was 0\% (200/200 passes) under the chosen generator/prompting policy and sampling semantics; (2) the adapter enforced shared\_initial semantics eliminated per‑condition sampling variance, so the absence of discordant pairs is not attributable to sampling differences across conditions but to the generator outputs themselves; and (3) the Mininet holdout found no validator--emulator mismatches on the 20‑task sample, providing limited end‑to‑end corroboration of validator pass verdicts for that holdout. All of these points are directly supported by the archived artifacts referenced above (execution/model\_configuration.json, deterministic\_reconciliation.json, analysis/condition\_summary.csv, analysis/paired\_contingency\_table.csv, analysis/contamination\_summary.csv, and representative task scoring artifacts).

\section{Discussion}
Conditional interpretation. The degenerate null outcome demonstrates a logical constraint: when the chosen generator, prompt template, and sampling policy produce validator‑passing candidates for every item in a workload, an added validator‑coupled repair step cannot reduce validator‑detected failures. This observation is artifact‑grounded and operationally important: the marginal value of automated repair depends on baseline validator failure prevalence and on validator coverage.

Why baseline failures were zero (artifact‑based explanations). First, the prompt template and enforced shared\_initial semantics produced outputs that conformed to explicit task‑bank constraints. Second, the task bank is curated with clear required\_sequences and machine‑checkable constraints that map directly to validator checks---this reduces semantic ambiguity and makes correct outputs easier to generate with a deterministic prompt policy. Third, the primary validator enforces the task‑encoded constraints but may omit device‑level idiosyncrasies detectable only by an end‑to‑end emulator or real device; the recorded Mininet holdout produced zero disagreements but the run also discloses a missing labeled selection seed for that holdout (see reproducibility note). Fourth, the single repair budget and shared\_initial design remove per‑condition re‑sampling as a pathway to obtain alternative repairable candidates.

Operational implications. Before deploying automated repair, practitioners should measure baseline validator failure prevalence on representative operational workloads and examine validator coverage with end‑to‑end emulation. When baseline validator failures are rare under the chosen generator/prompt policy, investing in richer validator fidelity (vendor semantics, device models), multi‑sample generation, multi‑model ensembles, or targeted adversarial testing may yield higher safety returns than automated repair. Conversely, when baseline failures are non‑negligible, a well‑engineered GVR pipeline with sufficient repair budget can reduce validator‑detected issues; realized gains will depend on repair‑prompt quality, budget, and validator completeness.

Threats to validity. Internal validity: preregistration, deterministic adapter enforcement, and frozen stopping rules minimize post‑hoc choices, but the shared\_initial design intentionally trades sampling generality for a controlled paired comparison. Construct validity: the primary validator enforces explicit task constraints but may not model all real device behaviors; emulator disagreements would reveal such blind spots but none were observed on the 20‑task Mininet holdout. External validity: a single hosted model (gpt‑5‑mini), single prompt/sampling policy (shared initial deterministic sampling), a curated synthetic/public task bank, and a single repair budget were used; results may differ with other models, nonzero temperature sampling, multiple samples per condition, adversarial workloads, or production device heterogeneity. Conclusion validity: with zero discordant pairs the McNemar test is uninformative for the preregistered effect size; nonetheless the observed zero‑discordance is a reproducible operational datum documented in the archived artifacts.

Reproducibility note and documented gap. All execution artifacts (model‑call logs, validator traces, paired contingency tables, bootstrap parameters) are archived in the execution bundle. The Mininet holdout evidence is captured in analysis/contamination\_summary.csv which reports zero disagreements for the 20‑task holdout; this artifact demonstrates that the secondary emulator check executed and found no validator--emulator mismatches. A remaining reproducibility gap is that the explicit randomized selection seed for the Mininet holdout is not labeled in the archived artifacts, preventing exact deterministic reconstruction of which tasks were drawn into the holdout despite the contamination summary reporting zero disagreements. We disclose this gap transparently in the Disclosure Statement.

\section{Limitations}
\begin{itemize}
\item Task‑bank scope: the confirmatory sample used a curated synthetic/public intent\ensuremath{\rightarrow}configuration bank (n = 200); results therefore reflect this bank and may not generalize to broader production workloads or adversarial distributions.
\item Single‑model and shared‑initial setting: only gpt‑5‑mini (hosted) with adapter‑enforced shared‑initial generation was tested (execution artifacts record hosted\_model\_call\_event\_count = 200 and stage\_counts shared\_initial\_generation = 200). Outcomes may differ across model families, independent per‑condition sampling, nonzero temperature sampling, or larger repair budgets.
\item Validator coverage and oracle limitations: the primary deterministic validator enforces a defined set of syntax/intent constraints; validators can fail to capture real device idiosyncrasies, vendor‑specific behaviors, or emergent failure modes detectable only by end‑to‑end deployment.
\item Adversarial robustness untested: this confirmatory run did not include active adversarial injection or red‑team tool‑description attacks; prior audits indicate such vectors can bypass naive defenses and remain a separate concern.
\item Reproducibility gap for contamination check: the randomized holdout selection seed for the Mininet contamination check is absent from the labeled artifacts, preventing exact reconstruction of the holdout selection.
\item Single repair budget: the GVR pipeline allowed at most one automated repair prompt per task; multi‑step repair strategies, different repair budgets, or human‑in‑the‑loop approvals were not exercised here.
\end{itemize}

\section{Conclusion}
We preregistered and executed a sealed paired study comparing LLM‑only generation to a deterministic GVR pipeline on 200 NetOps intent\ensuremath{\rightarrow}configuration tasks using gpt‑5‑mini and deterministic\_netops\_validator\_v5. Both pipelines produced validator‑passing configurations for every task (n11 = 200; no discordant pairs), resulting in an observed paired difference of 0.0, exact McNemar p = 1.0, and bootstrap 95\% CI = [0.0, 0.0]. The preregistered >=50\% relative‑reduction hypothesis could not be evaluated because baseline validator failures were absent. The artifact‑grounded outcome clarifies that GVR's marginal value is workload dependent: when baseline validator‑detectable failures are rare, automated repair yields no measured benefit for those validator‑detected errors. Future studies should target workloads with nontrivial baseline failure rates, expand validator fidelity with richer emulation or real‑device checks, evaluate multi‑sample and multi‑model policies, and explore multi‑step or human‑in‑the‑loop repair budgets to characterize practical operational gains. All execution artifacts and analysis outputs are archived for independent inspection under the preregistration id cited above.

\section*{Disclosure Statement}
Disclosure (mandatory). This study executed under the autonomous post‑lock preregistration cnsms2026‑gvr‑effectiveness‑02. Declared model: gpt‑5‑mini (provider label: openai\_responses) invoked via the hosted\_netops\_gvr\_v1 adapter. The exact initial master prompt is archived at provenance/master\_prompt.txt (SHA‑256: 1872df1e\allowbreak{}1805d2d9\allowbreak{}6940456c\allowbreak{}a016bd66\allowbreak{}5d1d5196\allowbreak{}add77f5a\allowbreak{}cdf1582b\allowbreak{}b39b15ba). Topic constraints: the human Research Director limited the study to Generative AI and LLMs for NetOps focusing on safety/reliability of generated network changes. Frameworks and tools: orchestration used hosted\_netops\_gvr\_v1 and the deterministic task generator deterministic\_netops\_task\_generator\_v5. Primary deterministic validator: deterministic\_netops\_validator\_v5 (intent/constraint checks). A Mininet‑based emulator was preregistered and executed as a secondary disagreement detector on a randomized 10\% holdout (20 tasks); the artifact analysis/contamination\_summary.csv shows zero disagreements for that holdout. Preregistration and freeze: the experimental plan (estimand, sample size n = 200, paired design, stopping rules, max model calls = 400, repair budget = 1 per task, shared‑initial generation semantics) was frozen before any model calls. Autonomous execution and artifacts: the run executed autonomously under the frozen plan (start: 2026‑08‑30T11:50:34.530059Z; end: 2026‑08‑30T12:12:07.541650Z); model calls, prompts, seeds, validator traces, emulator traces, logs, and the artifact manifest are archived in the execution bundle. Model‑call budget and usage: 200 model calls were used (one shared initial generation per task); up to 200 repair calls were allowed but none were applied because initial candidates passed the validator. Human involvement: the human Research Director selected the topic family and pre‑lock constraints; no human manual editing or selective rewriting of technical manuscript content occurred after the autonomy lock. Deviations and restarts: no post‑preregistration design changes or technical restarts occurred. Known reproducibility gap: the explicit randomized selection seed used to draw the Mininet 10\% holdout is not present as a labeled artifact in the archive; this absence prevents exact deterministic reconstruction of the drawn holdout despite the contamination summary reporting no disagreements. The master prompt archive reference above is the immutable prompt required by the track.

\begin{thebibliography}{99}
\bibitem{ref1} Wayne Xin Zhao, Kun Zhou, Junyi Li, Tianyi Tang, Zican Dong, Yupeng Hou, Beichen Zhang, Yingqian Min, Junjie Zhang, Peiyu Liu, Xiaolei Wang, Yifan Du, Yushuo Chen, Yushuo Chen, Zhipeng Chen, Jinhao Jiang, Ruiyang Ren, Yifan Li, Xinyu Tang, Peiyu Liu, Yiwen Hu, Jian‐Yun Nie, Ji-Rong Wen, "A Survey of Large Language Models", 2026, 10.1007/s11704-026-60308-3
\bibitem{ref2} Changjie Wang, Mariano Scazzariello, Alireza Farshin, Simone Ferlin, Dejan Kostić, Marco Chiesa, "NetConfEval: Can LLMs Facilitate Network Configuration?", 2024, 10.1145/3656296
\bibitem{ref3} Oscar G. Lira, Oscar Maurício Caicedo Rendón, Nelson L. S. da Fonseca, "Large Language Models for Zero Touch Network Configuration Management", 2024, 10.1109/mcom.001.2400368
\bibitem{ref4} Chirag Shah, Ryen W. White, Reid Andersen, Georg Buscher, Scott Counts, Sarkar Snigdha Sarathi Das, Ali Montazeralghaem, Sathish Manivannan, J. Neville, Nagu Rangan, Tara Safavi, Siddharth Suri, Mengting Wan, Leijie Wang, Longqi Yang, "Using Large Language Models to Generate, Validate, and Apply User Intent Taxonomies", 2025, 10.1145/3732294
\bibitem{ref5} omar altawil, Dr. Mustafa Al-Fayoumi, "Configuration-Level Semantic Brittleness in Tool-Using LLM Agents: A Factor-Ranking Study of Persona, History, and Tool Definition", 2026, 10.2139/ssrn.7203631
\bibitem{ref6} Mohammadreza Rashidi, "Measuring How LLM Tool Descriptions, Cross-Tool Ambiguity, and Action Chains Compose into Excessive Agency Exploits", 2026, 10.21203/rs.3.rs-10646209/v1
\bibitem{ref7} Yoshihiro Ando, "Secure-by-Default Guardrails for MCP-Based Tool Use in Multi-Modal LLM Agents", 2026, 10.36227/techrxiv.177006109.97957916/v1
\bibitem{ref8} Sebastian Farquhar, Jannik Kossen, Lorenz Kuhn, Yarin Gal, "Detecting hallucinations in large language models using semantic entropy", 2024, 10.1038/s41586-024-07421-0
\bibitem{ref9} Shuyin Ouyang, Jie M. Zhang, Mark Harman, Meng Wang, "An Empirical Study of the Non-Determinism of ChatGPT in Code Generation", 2024, 10.1145/3697010
\bibitem{ref10} Oghenekeno Hilkiah Okula, ""The Era of AI Agents for Network Engineering": Intent-Aware Auto Remediation for Network Disruption or Failures Using AI Agents, Large Language Models (LLM) and Model Context Protocol (MCP) Mechanisms", 2026, 10.36227/techrxiv.177004277.71795055/v1
\end{thebibliography}

\end{document}
